\section{注释}

\begin{enumerate}
\item\hypertarget{endnote:1}{} J.L. Collins said it best in a remake of this classic scene from the film The Gambler: www.youtube.com/watch?v=eikbQPIdhPY
\item\hypertarget{endnote:2}{} Miller, Matthew L., “Binge 'Til You Burst: Feast and Famine on Salmon Rivers,” Cool Green Science (April 8, 2015).
\item\hypertarget{endnote:3}{} Nichols, Austin and Seth Zimmerman, “Measuring Trends in Income Variability,” Urban Institute Discussion Paper (2008).
\item\hypertarget{endnote:4}{} Dynan, Karen E., Jonathan Skinner, and Stephen P. Zeldes, “Do the Rich Save More?” Journal of Political Economy 112:2 (2004), 397–444.
\item\hypertarget{endnote:5}{} Saez, Emmanuel, and Gabriel Zucman, “The Distribution of US Wealth: Capital Income and Returns since 1913.” Unpublished (2014).
\item\hypertarget{endnote:6}{} “Stress in America? Paying With Our Health,” American Psychological Association (February 4, 2015).
\item\hypertarget{endnote:7}{} “Planning \& Progress Study 2018,” Northwestern Mutual (2018).
\item\hypertarget{endnote:8}{} Graham, Carol, “The Rich Even Have a Better Kind of Stress than the Poor,” Brookings (October 26, 2016).
\item\hypertarget{endnote:9}{} Leonhardt, Megan, “Here's How Much Money Americans Say You Need to Be 'Rich',” CNBC (July 19, 2019).
\item\hypertarget{endnote:10}{} Frank, Robert, “Millionaires Need \$7.5 Million to Feel Wealthy,” The Wall Street Journal (March 14, 2011).
\item\hypertarget{endnote:11}{} Chris Browning et al., “Spending in Retirement: Determining the Consumption Gap,” Journal of Financial Planning 29:2 (2016), 42.
\item\hypertarget{endnote:12}{} Taylor, T., Halen, N., and Huang, D., “The Decumulation Paradox: Why Are Retirees Not Spending More?” Investments \& Wealth Monitor (July/August 2018), 40-52.
\item\hypertarget{endnote:13}{} Matt Fellowes, “Living Too Frugally? Economic Sentiment \& Spending Among Older Americans,” unitedincome.capitalone.com (May 16, 2017).
\item\hypertarget{endnote:14}{} Survey of Consumer Finances and Financial Accounts of the United States.
\item\hypertarget{endnote:15}{} 19th Annual Transamerica Retirement Survey (December 2019).
\item\hypertarget{endnote:16}{} The 2020 Annual Report of the Board of Trustees of the Federal Old-Age and Survivors Insurance and Federal Disability Insurance Trust Funds (April 2020).
\item\hypertarget{endnote:17}{} Pontzer, Herman, David A. Raichlen, Brian M. Wood, Audax Z.P. Mabulla, Susan B. Racette, and Frank W. Marlowe, “Hunter-gatherer Energetics and Human Obesity,” PLoS ONE 7:7 (2012), e40503.
\item\hypertarget{endnote:18}{} Ross, Robert, and I.N. Janssen, “Physical Activity, Total and Regional Obesity: Dose-response Considerations,” Medicine and Science in Sports and Exercise 33:6 SUPP (2001), S521-S527.
\item\hypertarget{endnote:19}{} Balboni, Clare, Oriana Bandiera, Robin Burgess, Maitreesh Ghatak, and Anton Heil, “Why Do People Stay Poor?” (2020). CEPR Discussion Paper No. DP14534.
\item\hypertarget{endnote:20}{} Egger, Dennis, Johannes Haushofer, Edward Miguel, Paul Niehaus, and Michael W. Walker, “General Equilibrium Effects of Cash Transfers: Experimental Evidence From Kenya,” No. w26600. National Bureau of Economic Research (2019).
\item\hypertarget{endnote:21}{} Stanley, Thomas J., The Millionaire Next Door: The Surprising Secrets of America's Wealthy (Lanham, MD: Taylor Trade Publishing, 1996).
\item\hypertarget{endnote:22}{} Corley, Thomas C., “It Takes the Typical Self-Made Millionaire at Least 32 Years to Get Rich,” Business Insider (March 5, 2015).
\item\hypertarget{endnote:23}{} Curtin, Melanie, “Attention, Millennials: The Average Entrepreneur is This Old When They Launch Their First Startup,” Inc.com (May 17, 2018).
\item\hypertarget{endnote:24}{} Martin, Emmie, “Suze Orman: If You Waste Money on Coffee, It's like 'Peeing \$1 Million down the Drain',” CNBC (March 28, 2019).
\item\hypertarget{endnote:25}{} Rigby, Rhymer, “We All Have Worries but Those of the Rich Are Somehow Different,” Financial Times (February 26, 2019).
\item\hypertarget{endnote:26}{} Dunn, Elizabeth, and Michael I. Norton, Happy Money: The Science of Happier Spending (New York, NY: Simon \& Schuster Paperbacks, 2014).
\item\hypertarget{endnote:27}{} Pink, Daniel H, Drive: The Surprising Truth about What Motivates Us (New York, NY: Riverhead Books, 2011).
\item\hypertarget{endnote:28}{} Matz, Sandra C., Joe J. Gladstone, and David Stillwell, “Money Buys Happiness When Spending Fits Our Personality, Psychological Science 27:5 (2016), 715-725.
\item\hypertarget{endnote:29}{} Dunn, Elizabeth W., Daniel T. Gilbert, and Timothy D. Wilson, “If Money Doesn’t Make You Happy, Then You Probably Aren’t Spending It Right,” Journal of Consumer Psychology 21:2 (2011), 115–125.
\item\hypertarget{endnote:30}{} Vanderbilt, Arthur T., Fortune’s Children: The Fall of the House of Vanderbilt (New York, NY: Morrow, 1989).
\item\hypertarget{endnote:31}{} Gorbachev, Olga, and María José Luengo-Prado, “The Credit Card Debt Puzzle: The Role of Preferences, Credit Access Risk, and Financial Literacy,” Review of Economics and Statistics 101:2 (2019), 294-309.
\item\hypertarget{endnote:32}{} Collins, Daryl, Jonathan Morduch, Stuart Rutherford, and Orlanda Ruthven, Portfolios of the Poor: How the World’s Poor Live on \$2 a Day (Princeton, NJ: Princeton University Press, 2009).
\item\hypertarget{endnote:33}{} “The Economic Value of College Majors,” CEW Georgetown (2015).
\item\hypertarget{endnote:34}{} Tamborini, Christopher R., ChangHwan Kim, and Arthur Sakamoto, “Education and Lifetime Earnings in the United States,” Demography 52:4 (2015), 1383-1407.
\item\hypertarget{endnote:35}{} “The Economic Value of College Majors,” CEW Georgetown (2015).
\item\hypertarget{endnote:36}{} “Student Loan Debt Statistics [2021]: Average + Total Debt,” EducationData (April 12, 2020).
\item\hypertarget{endnote:37}{} Radwin, David, and C. Wei, “What is the Price of College? Total, Net, and Out-of-Pocket Prices by Type of Institution in 2011-12,” Resource document, National Center for Education Statistics (2015).
\item\hypertarget{endnote:38}{} Brown, Sarah, Karl Taylor, and Stephen Wheatley Price, “Debt and Distress: Evaluating the Psychological Cost of Credit,” Journal of Economic Psychology 26:5 (2005), 642-663.
\item\hypertarget{endnote:39}{} Dunn, Lucia F., and Ida A. Mirzaie, “Determinants of Consumer Debt Stress: Differences by Debt Type and Gender,” Department of Economics: Columbus, Ohio State University (2012).
\item\hypertarget{endnote:40}{} Sweet, Elizabeth, Arijit Nandi, Emma K. Adam, and Thomas W. McDade, “The High Price of Debt: Household Financial Debt and its Impact on Mental and Physical Health,” Social Science \& Medicine 91 (2013), 94-100.
\item\hypertarget{endnote:41}{} Norvilitis, J.M., Szablicki, P.B., and Wilson, S.D., “Factors Influencing Levels of Credit-Card Debt in College Students,” Journal of Applied Social Psychology 33 (2003), 935-947.
\item\hypertarget{endnote:42}{} Dixon, Amanda, “Survey: Nearly 4 in 10 Americans Would Borrow to Cover a \$1K Emergency,” Bankrate (January 22, 2020).
\item\hypertarget{endnote:43}{} Kirkham, Elyssa, “Most Americans Can't Cover a \$1,000 Emergency With Savings,” LendingTree (December 19, 2018).
\item\hypertarget{endnote:44}{} Athreya, Kartik, José Mustre-del-Río, and Juan M. Sánchez, “The Persistence of Financial Distress,” The Review of Financial Studies 32:10 (2019), 3851-3883.
\item\hypertarget{endnote:45}{} Shiller, Robert J., “Why Land and Homes Actually Tend to Be Disappointing Investments,” The New York Times (July 15, 2016).
\item\hypertarget{endnote:46}{} Bhutta, Neil, Jesse Bricker, Andrew C. Chang, Lisa J. Dettling, Sarena Goodman, Joanne W. Hsu, Kevin B. Moore, Sarah Reber, Alice Henriques Volz, and Richard Windle, “Changes in US Family Finances from 2016 to 2019: Evidence From the Survey of Consumer Finances,” Federal Reserve Bulletin 106:5 (2020).
\item\hypertarget{endnote:47}{} Eggleston, Jonathan, Donald Hayes, Robert Munk, and Brianna Sullivan, “The Wealth of Households: 2017,” U.S. Census Bureau Report P70BR-170 (2020).
\item\hypertarget{endnote:48}{} Kushi, Odeta, “Homeownership Remains Strongly Linked to Wealth-Building,” First American (November 5, 2020).
\item\hypertarget{endnote:49}{} “What Is a Debt-to-Income Ratio? Why Is the 43\% Debt-to-Income Ratio Important?” Consumer Financial Protection Bureau (November 15, 2019).
\item\hypertarget{endnote:50}{} Bengen, W. P., “Determining Withdrawal Rates Using Historical Data,” Journal of Financial Planning 7:4 (1994), 171–182.
\item\hypertarget{endnote:51}{} Kitces, Michael, “Why Most Retirees Never Spend Their Retirement Assets,” Nerd's Eye View, Kitces.com (July 6, 2016).
\item\hypertarget{endnote:52}{} Bengen, William, Interview with Michael Kitces, Financial Advisor Success Podcast (October 13, 2020).
\item\hypertarget{endnote:53}{} “Spending in Retirement,” J.P. Morgan Asset Management (August 2015).
\item\hypertarget{endnote:54}{} Fisher, Jonathan D., David S. Johnson, Joseph Marchand, Timothy M. Smeeding, and Barbara Boyle Torrey, “The Retirement Consumption Conundrum: Evidence From a Consumption Survey,” Economics Letters 99:3 (2008), 482-485.
\item\hypertarget{endnote:55}{} Robin, Vicki, Joe Dominguez, and Monique Tilford, Your Money or Your Life: 9 Steps to Transforming Your Relationship with Money and Achieving Financial Independence (Harmondsworth: Penguin, 2008).
\item\hypertarget{endnote:56}{} Zelinski, Ernie J., How to Retire Happy, Wild, and Free: Retirement Wisdom That You Won't (Visions International Publishing: 2004).
\item\hypertarget{endnote:57}{} O'Leary, Kevin, “Kevin O'Leary: Why Early Retirement Doesn't Work,” YouTube video, 1:11 (March 20, 2019).
\item\hypertarget{endnote:58}{} Shapiro, Julian, “Personal Values,” Julian.com.
\item\hypertarget{endnote:59}{} Maggiulli, Nick, “If You Play With FIRE, Don't Get Burned.” Of Dollars And Data (March 30, 2021).
\item\hypertarget{endnote:60}{} “Social Security Administration,” Social Security History, ssa.gov.
\item\hypertarget{endnote:61}{} Roser, M., Ortiz-Ospina, E., and Ritchie, H., “Life Expectancy,” ourworldindata.org (2013).
\item\hypertarget{endnote:62}{} Hershfield, Hal E., Daniel G. Goldstein, William F. Sharpe, Jesse Fox, Leo Yeykelis, Laura L. Carstensen, and Jeremy N. Bailenson, “Increasing Saving Behavior Through Age- Progressed Renderings of the Future Self,” Journal of Marketing Research 48 SPL (2011), S23-S37.
\item\hypertarget{endnote:63}{} Fisher, Patti J., and Sophia Anong, “Relationship of Saving Motives to Saving Habits,” Journal of Financial Counseling and Planning 23:1 (2012).
\item\hypertarget{endnote:64}{} Colberg, Fran, “The Making of a Champion,” Black Belt (April 1975).
\item\hypertarget{endnote:65}{} Seigel, Jeremy J., Stocks for the Long Run (New York, NY: McGraw-Hill, 2020).
\item\hypertarget{endnote:66}{} Dimson, Elroy, Paul Marsh, and Mike Staunton, Triumph of the Optimists: 101 Years of Global Investment Returns (Princeton, NJ: Princeton University Press, 2009).
\item\hypertarget{endnote:67}{} Biggs, Barton, Wealth, War and Wisdom (Oxford: John Wiley \& Sons, 2009).
\item\hypertarget{endnote:68}{} U.S. Department of the Treasury, Daily Treasury Yield Curve Rates (February 12, 2021).
\item\hypertarget{endnote:69}{} Asness, Clifford S., “My Top 10 Peeves,” Financial Analysts Journal 70:1 (2014), 22–30.
\item\hypertarget{endnote:70}{} Jay Girotto, interview with Ted Seides, Capital Allocators, podcast audio (October 13, 2019).
\item\hypertarget{endnote:71}{} Beshore, Brent (@brentbeshore). 12 Dec 2018, 3:52 PM. Tweet.
\item\hypertarget{endnote:72}{} Wiltbank, Robert, and Warren Boeker, “Returns To Angel Investors In Groups,” SSRN.com (November 1, 2007); and “Review of Research on the Historical Returns of the US Angel Market,” Right Side Capital Management, LLC (2010).
\item\hypertarget{endnote:73}{} “Who Are American Angels? Wharton and Angel Capital Association Study Changes Perceptions About the Investors Behind U.S. Startup Economy,” Angel Capital Association (November 27, 2017).
\item\hypertarget{endnote:74}{} Altman, Sam, “Upside Risk,” SamAltman.com (March 25, 2013).
\item\hypertarget{endnote:75}{} Max, Tucker, “Why I Stopped Angel Investing (and You Should Never Start),” Observer.com (August 11, 2015).
\item\hypertarget{endnote:76}{} Wiltbank, Robert, and Warren Boeker, “Returns To Angel Investors in Groups,” SSRN.com (November 1, 2007).
\item\hypertarget{endnote:77}{} Frankl-Duval, Mischa, and Lucy Harley-McKeown, “Investors in Search of Yield Turn to Music-Royalty Funds,” The Wall Street Journal (September 22, 2019).
\item\hypertarget{endnote:78}{} SPIVA, spglobal.com (June 30, 2020).
\item\hypertarget{endnote:79}{} Bessembinder, Hendrik, “Do Stocks Outperform Treasury Bills?” Journal of Financial Economics 129:3 (2018), 440–457.
\item\hypertarget{endnote:80}{} West, Geoffrey B., Scale: The Universal Laws of Life, Growth, and Death in Organisms, Cities, and Companies (Harmondsworth: Penguin, 2017).
\item\hypertarget{endnote:81}{} Kosowski, Robert, Allan Timmermann, Russ Wermers, and Hal White, “Can Mutual Fund 'Stars' Really Pick Stocks? New Evidence from a Bootstrap Analysis,” The Journal of Finance 61:6 (2006), 2551-2595.
\item\hypertarget{endnote:82}{} “The Truth About Top-Performing Money Managers,” Baird Asset Management, White Paper (2014).
\item\hypertarget{endnote:83}{} Powell, R., “Bernstein: Free Trading is Like Giving Chainsaws to Toddlers.” The Evidence-Based Investor (March 25, 2021).
\item\hypertarget{endnote:84}{} Stephens-Davidowitz, Seth, Everybody Lies: Big Data, New Data, and What the Internet Can Tell Us About Who We Really Are (New York, NY: HarperCollins, 2017).
\item\hypertarget{endnote:85}{} Rosling, Hans, Factfulness (Paris: Flammarion, 2019).
\item\hypertarget{endnote:86}{} Buffett, Warren E., “Buy American. I Am,” The New York Times (October 16, 2008).
\item\hypertarget{endnote:87}{} “Asset Allocation Survey,” aaii.com (March 13, 2021).
\item\hypertarget{endnote:88}{} This is the median outcome for investing every month for a decade into U.S. stocks from 1926-2020.
\item\hypertarget{endnote:89}{} For more detail, see: ofdollarsanddata.com/in-defense-of-global-stocks.
\item\hypertarget{endnote:90}{} Zax, David, “How Did Computers Uncover J. K. Rowling’s Pseudonym?” Smithsonian Institution, Smithsonian.com (March 1, 2014).
\item\hypertarget{endnote:91}{} Hern, Alex, “Sales of ‘The Cuckoo’s Calling’ Surge by 150,000\% After J. K. Rowling Revealed as Author,” New Statesman (July 14, 2013).
\item\hypertarget{endnote:92}{} Kitces, Michael, “Understanding Sequence of Return Risk \& Safe Withdrawal Rates,” Kitces.com (October 1, 2014).
\item\hypertarget{endnote:93}{} Frock, Roger, Changing How the World Does Business: FedEx's Incredible Journey to Success - The Inside Story (Oakland, CA: Berrett-Koehler Publishers, 2006).
\item\hypertarget{endnote:94}{} Anarkulova, Aizhan, Scott Cederburg, and Michael S. O’Doherty, “Stocks for the Long Run? Evidence from a Broad Sample of Developed Markets,” SSRN.com (May 6, 2020).
\item\hypertarget{endnote:95}{} Zilbering, Yan, Colleen M. Jaconetti, and Francis M. Kinniry Jr., “Best Practices for Portfolio Rebalancing,” Valley Forge, PA: The Vanguard Group.
\item\hypertarget{endnote:96}{} Bernstein, William J., “The Rebalancing Bonus,” www. efficientfrontier.com.
\item\hypertarget{endnote:97}{} Brownlee, W. Elliot, Federal Taxation in America (Cambridge: Cambridge University Press, 2016).
\item\hypertarget{endnote:98}{} Leonhardt, Megan, “Here's What the Average American Typically Pays in 401(k) Fees,” CNBC (July 22, 2019).
\item\hypertarget{endnote:99}{} Witt, April, “He Won Powerball’s \$314 Million Jackpot. It Ruined His Life,” The Washington Post (October 23, 2018).
\item\hypertarget{endnote:100}{} Luce, Edward, “Lloyd Blankfein: I Might Find It Harder to Vote for Bernie than for Trump,” Financial Times (February 21, 2020).
\item\hypertarget{endnote:101}{} Saez, Emmanuel, and Gabriel Zucman, “Wealth Inequality in the United States Since 1913: Evidence from Capitalized Income Tax Data, 1913–2013,” The Quarterly Journal of Economics 131:2 (2016), 519–578.
\item\hypertarget{endnote:102}{} Karadja, Mounir, Johanna Mollerstrom, and David Seim, “Richer (and Holier) Than Thou? The Effect of Relative Income Improvements on Demand for Redistribution,” Review of Economics and Statistics 99:2 (2017), 201-212.
\item\hypertarget{endnote:103}{} Jackson, Matthew O., The Human Network: How Your Social Position Determines Your Power, Beliefs, and Behaviors (New York, NY: Vintage, 2019).
\item\hypertarget{endnote:104}{} “Global Wealth Report 2018.” Credit Suisse (October 18, 2018).
\item\hypertarget{endnote:105}{} Peter Attia, “Reverse Engineered Approach to Human Longevity,” YouTube video, 1:15:37 (November 25, 2017).
\item\hypertarget{endnote:106}{} Guvenen, Fatih, Fatih Karahan, Serdar Ozkan, and Jae Song, “What Do Data on Millions of US Workers Reveal About Life-cycle Earnings Dynamics?” FRB of New York Staff Report 710 (2015).
\item\hypertarget{endnote:107}{} Schwandt, Hannes, “Human Wellbeing Follows a U-Shape over Age, and Unmet Aspirations Are the Cause,” British Politics and Policy at LSE (August 7, 2013).
\item\hypertarget{endnote:108}{} Rauch, Jonathan, The Happiness Curve: Why Life Gets Better After 50 (New York, NY: Thomas Dunne Books, 2018).
\end{enumerate}
